%-------------------------------------------------------------------------------
%\tSECTION TITLE
%-------------------------------------------------------------------------------
\cvsection{Expérience Professionnelle}

%-------------------------------------------------------------------------------
%   CONTENT
%-------------------------------------------------------------------------------
\begin{cventries}
  %---------------------------------------------------------
  \cventry
  {Service d'appui à la recherche - Laboratoire d'Informatique et de Robotique de l'Université de
  Montpellier} % Organization
  {Ingénieur de Recherche (CDD -- BIATSS)} % Job title
  {Montpellier, France} % Location
  {Fév. 2024 -- Aujourd'hui} % Date(s)
  {
    \begin{cvitems} % Description(s) of tasks/responsibilities
    \item \entrypositionstyle{Développement Logiciel} : \\
      -- Mainteneur principal de l'écosystème de contrôle robotique temps-réel \keyword{mc\_rtc} permettant
      le contrôle de plusieurs dizaines de robots (humanoïdes, manipulateurs, etc) et activement utilisé dans
      plusieurs laboratoires (\keyword{LIRMM, JRL, EPFL, TUM}, ...) et pour des projets européens d'envergure
      (\keyword{H2020 COMANOID}, \keyword{I.AM}, etc).\\
      -- Co-conception avec le LAAS d'une architecture moderne de développement et de déploiement
      reproductible (\keyword{Nix}).
    \item \entrypositionstyle{Plateformes} : \\
      -- Responsable des robots humanoïdes (\keyword{HRP-4, RHPS1, Unitree G1, Pepper, NAO}) et des bras
      manipulateurs (\keyword{Panda, UR5}). \\
      -- Effectue l'entretien mécanique et électronique régulier. \\
      -- Effectue des réparations mécaniques et électroniques simples: remplacement de pièces (e.g moteur),
      remplacement de câbles électroniques défectueux, etc. \\
      -- Assure le maintien logiciel de ces systèmes complexes pour garantir leur disponibilité immédiate
      lors des expérimentations.
    \item \entrypositionstyle{Encadrement} : Accompagnement de doctorants pour la
      mise en œuvre de démonstrateurs expérimentaux :\\
      -- Projet d'assistance à se lever d'une chaise. Garant des enjeux de sécurité lors des interactions
      robot-humain. \\
      -- Mise en place d'un démonstrateur de stabilité dynamique sur \emph{HRP-4} et la nouvelle plateforme
      \emph{RHPS1}. Valorisation lors de l'évaluation HCERES du laboratoire.
    \item \entrypositionstyle{Communication} - Présentations techniques à :\\
      -- \emph{Réseau métier des ingénieurs roboticiens - Tech Days 2RM} - \keyword{Theme}: Reproducibilité
      logicielle. Excellente réceptivité de la communauté. \\
      -- \emph{ROSCon Strasbourg 2025}: \keyword{Theme}: contrôle par optimisation quadratique (QP) dans ROS2.
    \item \entrypositionstyle{Projets}:
      \begin{itemize}
        \item \entrypositionstyle{ANR Rolkneematics} : \\
          -- Mise en place d'un banc d'essai de suivi de trajectoires de genou. \\
          -- Acquisition d'un dataset de données capteur (partenariat BoneTag) et données de trajectories
          pour apprentissage IA. \\
          -- Co-ordination de 2 ingénieurs CDD (mécatronique, IA) et un postdoc (modélisation magnétique).
      \end{itemize}
    \end{cvitems}
  }
  %---------------------------------------------------------
  \cventry
  {CNRS -- LIRMM -- Équipe Interactive Digital Human (IDH)} % Organization
  {Ingénieur de Recherche (CDD)} % Job title
  {Montpellier, France} % Location
  {Nov. 2021 -- Avr. 2023} % Date(s)
  {
    \begin{cvitems} % Description(s) of tasks/responsibilities
    \item \entrypositionstyle{Développement Logiciel} :
      \begin{subitems}
      \item Co-maintenance et développement actif de \keyword{mc\_rtc}, support technique auprès des équipes du LIRMM.
      \end{subitems}
    \item \entrypositionstyle{Encadrement} : Accompagnement technologique sur 5 projets doctoraux, dont :
      \begin{subitems}
      \item Développement d'un banc d'essai robotique pour tester et valider la conception mécanique d'une orthèse.
      \item Suivi et déploiement d'un système de téléopération du robot \keyword{Pepper} pour une étude
        d'interaction avec enfants autistes au \emph{CHU de Montpellier}.
      \end{subitems}
    \item \entrypositionstyle{Projets} - Responsable de l'intégration robotique pour :
      \begin{subitems}
      \item Deux projets en collaboration avec \keyword{Michelin} avec démonstrateurs finaux validés sur le
        site de Michelin Ladoux:
        \begin{subitems}
        \item Projet de déchargement automatisé de cargaisons de pneumatiques en collaboration avec
          \keyword{Michelin}. Démonstration sur site approuvée par \keyword{Michelin}.
        \item Démonstrateur final sur site industriel de \keyword{Michelin Ladoux} d'un projet de
          loco-manipulation de gros objets industriels (bobine de 110Kg) et d'une nappe de caoutcouc flexible.
        \end{subitems}
      \item Une mise-en-scène de danse interactive humain-robot pour la soirée d'ouverture saisonnière 2024
        du Centre des Arts d'Enghien-les-Bains.
      \end{subitems}
    \end{cvitems}
  }

  %---------------------------------------------------------
  \cventry
  {AIST -- Joint Robotics Laboratory (JRL)} % Organization
  {Ingénieur de Recherche (CDD)} % Job title
  {Tsukuba, Japon} % Location
  {Nov. 2019 -- Avr. 2021} % Date(s)
  {
    \begin{cvitems} % Description(s) of tasks/responsibilities
    \item \entrypositionstyle{Développement Logiciel} : \\
      -- Pilotage stratégique de l'unification du framework de contrôle historique de l'AIST (\emph{HMC})
      avec celui du JRL (\emph{mc\_rtc}) afin de rationaliser les développements. \\
      -- Portage d'un démonstrateur clé de vissage industriel sur \emph{HRP-5P} (asservissement visuel,
      contrôle hybride force/position, étalonnage des capteurs). \\
      -- Pérennisation de la brique logicielle d'asservissement visuel.
    \item \entrypositionstyle{Plateformes} : Intégration du nouveau robot humanoïde \emph{HRP-5P} au
      framework de contrôle.
    \item \entrypositionstyle{Projets} : \\
      -- \emph{ANA Avatar X-Prize (Equipe Janus)} : Concepteur principal de l'architecture de téléopération
      du robot \emph{HRP-4CR}. Sélectionnés en demi-finale. \\
      -- \emph{Michelin - loco-manipulation} : Responsable de l'intégration logicielle pour un projet de
      loco-manipulation d'objets lourds. Démonstrateur au laboratoire validé par \keyword{Michelin}.
    \end{cvitems}
  }

  %---------------------------------------------------------
  \cventry
  {CNRS -- LIRMM -- Équipe Interactive Digital Human (IDH)} % Organization
  {Ingénieur de Recherche (CDD)} % Job title
  {Montpellier, France} % Location
  {Oct. 2018 -- Oct. 2019} % Date(s)
  {
    \begin{cvitems} % Description(s) of tasks/responsibilities
    \item \entrypositionstyle{Développement Logiciel} : Premières contributions au framework
      \keyword{mc\_rtc}: outils de
      visualisation 3D temps réel et système d'observation d'état.
    \item \entrypositionstyle{Encadrement} : Mise en place de 3 démonstrateurs physiques sur le robot
      humanoïde \emph{HRP-4} pour appuyer et valider des publications scientifiques : \\
      -- contrôle des forces de contact en glissement (prédiction, adaptation dynamique du centre de masse)\\
      -- nouvelle formulation de contrainte d'évitement de collision et sa validation au contrôle en boucle fermée \\
      -- démonstrateur industriel de pick-and-drop logistique incluant l'exploitation de l'impact
      robot-environnement (contraintes QP des forces maximales post-impact admissibles).
    \item \entrypositionstyle{Projets} : \\
      -- Rôle clé dans le projet Européen \keyword{H2020 COMANOID} avec l'intégration de mes travaux de thèse
      pour assurer la localisation fiable du robot (SLAM) lors du démonstrateur final du projet sur le site
      industriel d'Airbus Saint-Nazaire. Travaux d'équipe récompensés par le prix du meilleur article :
      \emph{IEEE Robotics and Automation Magazine’s 2019 Best
      Paper Award}.
    \end{cvitems}
  }
\end{cventries}
